% SOURCE_AUTHORITY: sga5_fr_workpass.tex, lines 7634--7659 (LF-normalized unit).
% SOURCE_SHA256: 791F4EFFC5E02832D5D77ED03518C8156D6F07E4C8238B03545DB93D883FBB28
\begin{prooftext}
Reduzimo-nos imediatamente ao caso em que \(B\) é o anel de polinômios \(A[T]\), dotado de sua graduação habitual. O \((B,\mathcal C)\)-módulo \(A[T]\otimes_A M=(M_n)_{n\in\mathbb N}\) é então obtido pondo \(M_n=M\) para todo \(n\), e fazendo com que \(T\) envie, para todo inteiro \(n\), \(M_n\) em \(M_{n+1}\) por meio da identidade de \(M\).

Seja \(N\) um sub-\((A[T],\mathcal C)\)-módulo graduado de \(A[T]\otimes_A M\). As componentes \((N_n)_{n\in\mathbb N}\) de \(N\) definem uma sucessão crescente
\[
N_0\subset N_1\subset\cdots\subset N_n\subset\cdots
\]
de subobjetos de \(M\). Suponhamos agora que \(M\) é noetheriano (resp. de tipo finito) e seja
\[
N^1\subset N^2\subset\cdots\subset N^p\subset\cdots
\]
uma sucessão crescente de sub-\((A[T],\mathcal C)\)-módulos graduados de \(A[T]\otimes_A M\) que, na primeira hipótese, quando \(M\) é de tipo finito, tem por união \(A[T]\otimes_A M\).

Consideremos então o quadro de entrada dupla que aparece a seguir:
\[
\begin{array}{cccccc}
N_0^1&\subset&N_0^2&\subset\cdots\subset&N_0^p&\subset\cdots\\
\cap&&\cap&&\cap\\
N_1^1&&\vdots&&\vdots\\
\vdots&&\vdots&&\vdots\\
\cap&&\cap&&\cap\\
N_n^1&\subset&N_n^2&\subset\cdots\subset&N_n^p&\subset\cdots .
\end{array}
\]
Pela hipótese sobre \(M\), é claro que existe um par \((p_0,n_0)\) tal que \(N^{p_0}_{n_0}\) seja maximal (resp. igual a \(M\)). Se \(p_1\) é um inteiro \(\ge p_0\) tal que os \(N_n^{p_1}\) (\(n\le n_0\)) sejam os maiores elementos das sucessões \((N_n^p)_{p\in\mathbb N}\), é claro que \(N^p=N^{p_1}\) sempre que \(p\ge p_1\), o que permite concluir em ambos os casos.
\end{prooftext}
